%%%%%%%%%%%%
% NOTES FOR GEOGECON CHAPTER
%%%%%%%%%%%%
\section{notes: geogecon \label{notes_ge}}

\subsection{The evolution of spatial economics: from thunen to the new economic geography}
\cite{fujita_evolution_2010}

Fujita uses the term `general location theory' as a catch-all for theories asking ``what types of economic activity take place where, and why?'' (p.1) The term actually comes from Isard, 1949. Fujita: ``A general location model aims, either in a descriptive or normative context, to explain the geographical distribution of all agents in a given location space, together with the associated spatial price system and trade pattern.'' (p.2)

Here's Fujita being explicit about this: ``I would like to claim that Thunen's theory is a \emph{theory of general equilibrium in space}, or for short, \emph{a general location theory}.'' (p.3)

This 'unified' approach is the location theory equivalent to general equilibrium: the aspiration is produce models able to say something about \emph{whole} systems. In his defence of the analytic approach to \gls{economic geography}, Overman neatly explains why this focus on `general' theories is different:

\begin{quote}
``Partial equilibrium reasoning (looking at just a small part of any problem) often fails to provide the right reasoning in a general equilibrium context (when taking the economy as a whole).'' \fixref{Overman, learn: not in bib.} 
\end{quote}

Most obviously - both in general equilibrium and location theory - the role of elasticities between goods is paramount. Considering only one good will entirely miss the impact of price changes in other goods. 

...

Fujita notes that, while location theory has a long history, general location approaches are still ``limited to several special classes of models''. (p.2) This is certainly true of \gls{ge}: as \citeauthor{fujita_spatial_2001} notes, \gls{ge} works with -

\begin{quote}
``... assumptions that reflect not so much a realistic view of how the world works as a judgement about what will make the analysis of geographic issues manageable without doing too much damage to the relevance of that analysis.'' (p.6)
\end{quote}

%section here on manageability, tractability. Perhaps Blaug paper on rigour vs relevance.
Mentions that GE ``embraces both increasing returns and imperfect competition, the two basic ingredients of the formation of actual economic landscapes.'' (2)

Characterises NEG as ``the most recent and most successful development of general location theory'', which I like. (3).

Worth noting p.8: von Thunen apparently pre-empting Marshall by some way on agglomerating forces.

makes a pretty hefty claim for GE: von Thunen could not have produced a general location theory, any more than one could expect ``de Vinci, upon one's first view of his amazing `flying machines' drawn in the fifteenth century, to invent a real airplane'' without access to the internal combustion engine. For Fujita, the Dixit-Stiglitz model ``provided location theory with a powerful combustion engine.'' (\citealt{fujita_evolution_2010} p.8)

Now, in my margin notes, I've put ``vs `games you can play with CES functions' '' - presumably to sarcastically draw attention to the contrast here. An avowedly narrow model - general, in the sense of 'whole world', but still narrow. Is the engine up to the job?

...

Notes that Weber ``did not consider any endogenous determination of the prices of inputs and outputs.'' (10) 

Big list of early C20th people working on theories of non-competitive firms - worth going back to. (10)

The line here: Hotelling's work on spatial competition makes Palander (1935) ``reconsider Weber's industrial location theory in the more general framework of spatial price theory.'' (11)

Fuj notes Palander found a different outcome to Hotelling for the two-seller problem - Palanders' results favour firms keeping a distance to avoid competition. Hotelling's findings were sequential, though, weren't they? Sellers always have the incentive to try and take a little more of the market. But there's clearly a trade-off. Hmm. 

Notes about Losch: ``no interdependence between markets of different goods was explicitly considered in terms of micro-economics.'' (13) No eeelasticiteee. ``Indeed even today, there is no satisfactory model of general location analysis with multiple oligopolistic industries.'' (13)

Same point as Isard - general equilibrium ``maintained the problem of space can be handled simply by defining each commodity by its location as well as by its physical characteristics, and hence once we have made a clever indexing of commodities, we can essentially forget space in economic theory.'' (14)

Fuj gives a great quote from Isard - Isard makes the argument about space ruling out perfect competition; 'respective position' has a huge effect on the level of competition. So, on top of needing some form of increasing returns, space really is a problem! (15)

A great quote again from Isard, where he identifies a 'general theory of monopolistic competition' as identical with a general theory of location.

...

Skipping merrily forward some way -




\subsection{Isard: The general theory of location and the space-economy}

So: the first attempt to define a 'general theory of location', which is sort of what I'm aiming for...

calls general equilibrium ``a wonderland of no dimensions'' (477) In discussing this, Isard notes that Hicks rejects monocomp models in favour of perfect comp, since the former ``introduces elements of indeterminacy, whereas his preference is for determinate solutions''. (Worth going back to Hicks to read the original. Hicks, value and capital, p.477) That's interesting because, of course, the ABM approach is most certainly not determinate in the analytic sense; it's probabalistic. 

He also notes a chat with Schumpeter: this is a reiteration of the argument that space costs are just costs, and can thus be considered along with other factors - there is nothing unique about them.

``In this sense the factor of space is repudiated, everything within the economy is in effect compressed to a point, and all spatial resistance disappears.'' (477) I've argued that geogecon does this too, though I note there's a section on 'continuous space' in 'the spatial economy' - I think that's 1D, racetrack...

Interesting: he even has a go at theories that look like ``a one point world, which somehow or other is conceived as divided into n parts, representing n nations, between which trade and trade barriers exist.'' (478) That's, of course, geogecon.

``We conceive the general theory of location and space- economy to be one which comprehends the economy in its totality. Not only are the mutual relations and interdependence of all economic elements, both in the aggregate and atomistically, of fundamental importance; but the spatial as well as the temporal (dynamic) character of the interrelated economic processes must enter the picture.'' (478)

Seems to be saying Chamberlin had a go at space, but treated it as ``one of the leading manifestations of the broader category of `product differentiation'.'' (479)

Interesting: Isard argues the first attempt at such a general theory is in ch.7 of Weber. Hmm... Isard ends this by saying ``The technique is inadequate; it does not present any general, heuristic principle by means of which one can order the spatial complexities involved in the total location of economic activities'' (481). That last sentence is good: `general, heuristic principle by means of which one can order the spatial complexities involved in the total location of economic activities.' That's what we're looking for.

Isard wants to find a ``common denomiinator in terms of which all the forces stemming from the various interrelations can be expressed and evaluated, and by means of which a net effect could perhaps be deduced.'' (482) My argument: well, micro-economics doesn't do bad at this. Some good ideas! Let's use those.

He does, however, support Weber's `evolutionary' approach, sequencing changes. I think that's actually pretty inescapable, and is what I'm looking at too. 

Hmm - gets more interesting. He discusses Engl\"{a}nder, and describing his work gives a great summary:

\begin{quote}
``Any given entrepreneur, in choosing the site at which to produce or render services, considers the various supply prices existing in the various localities for the inputs that he might possibly employ; at the same time he considers the various prices which might be obtained in the various localities for his product or services. When finally he does locate at a site, he influences in turn the prices of various inputs and outputs.''
\end{quote}

In a footnote he goes on to note: ``Further (1) by classifying raw materials and factors of production, whether mobile or immobile, as place-free (available everywhere under the same conditions), conditionally place-bound (available at all or some places, under unequal conditions) and unconditionally place-bound (present at one site), and (2) by conceiving immobile goods as goods of infinite weight which enter into production with infinite weight-loss, Englander brings together the specific location theories of industry and agriculture within the confines of his pure location theory, not as distinct compartments, but as internally related sectors.'' (483)

Some jolly good ideas there: infinite weight for immobile goods!

Notes now Pred\"{o}hl - ``wished to investigate how far the location problem is a price problem; location theory, a price theory. In other words, to what extent does the local distribution of production lie inside the economic relationship of interdependent prices?'' (483)

Which of course, it could be: but my point is, it's also a \emph{heterogeneity} theory. Space guarantees that. Which is to say, while it may be possible to reduce location to cost, the heterogeneity cannot be reduced without damaging model outcomes.












%END OF NOTES
